% TOP_PROBLEM: {"id": 1864, "title": "Gilbreath's Conjecture", "category_id": 1, "category": {"id": 1, "name": "number_theory", "display_name": "Number Theory", "description": "Properties of integers, prime numbers, Diophantine equations.", "slug": "number-theory", "order_index": 1, "created_at": "2026-07-31T15:26:25.670Z"}, "status": "open"}
% FIRST_POSTED: null
% ENTRY_KIND: statement_only
% Imported from: batch08.tex; UnsolvedMath ID 1864
\documentclass[11pt]{article}
\usepackage[margin=25mm]{geometry}
\usepackage{fontspec}
\setmainfont{Latin Modern Roman}
\usepackage{amsmath,amssymb,mathtools}
\usepackage{unicode-math}
\setmathfont{Latin Modern Math}
\usepackage{xurl,hyperref}
\hypersetup{colorlinks=true,urlcolor=blue}
\setcounter{secnumdepth}{0}
\setlength{\parindent}{0pt}
\setlength{\parskip}{5pt}
\setlength{\emergencystretch}{3em}
\newcommand{\sref}[2]{\hyperlink{source:#1:#2}{[S#2]}}
\newcommand{\cataloguescope}[1]{\paragraph{Recorded target.}#1}
\begin{document}
% UnsolvedMath ID 1864; source catalogue code GUY-A10
\setcounter{section}{1863}
\section{Gilbreath's Conjecture}\label{1864}
{\small\sffamily UnsolvedMath ID 1864 · Number Theory}
\subsection{Definitions and mathematical statement}

\paragraph{Shared notation.}$\mathbb N=\{1,2,\ldots\}$, and $\mathbb Z,\mathbb Q,\mathbb R,\mathbb C$ denote the integers, rationals, reals and complex numbers. Any other conventions are specified locally.


Let \(p_n\) be the \(n\)-th prime in increasing order
(\(p_1=2,p_2=3,\ldots\)). Form the infinite absolute-difference array
\[
 a_n^{(0)}=p_n,\qquad
 a_n^{(r+1)}=\bigl|a_{n+1}^{(r)}-a_n^{(r)}\bigr|
                  \quad(n\ge1,\ r\ge0).
\]
The zeroth row is the prime sequence, the first row is its consecutive
prime gaps, and each further row is the unsigned difference of the preceding row.
\paragraph{Statement recorded in the supplied source.}
Is
\[
                         a_1^{(r)}=1\qquad\text{for every }r\ge1?
\] \sref{1864}{2}
\subsection{Short English statement}
Starting with all primes in increasing order and repeatedly taking absolute differences, does every row after the initial prime row begin with one?
\subsection{Sources}
\begin{description}
\item[\hypertarget{source:1864:1}{S1}] ULAM.AI and the UnsolvedMath Contributors, \emph{UnsolvedMath: A Curated Collection of Open Mathematics Problems}, record 1864 (GUY-A10). CC BY 4.0. \url{https://huggingface.co/datasets/ulamai/UnsolvedMath}.
\item[\hypertarget{source:1864:2}{S2}] Zachary Chase, Zach Hunter and Terence Tao, \emph{Gilbreath's conjecture: a Cram\'er random model and a deterministic analysis}, arXiv:2607.08712 (2026), Section 1.1, absolute-difference recursion (1.1)--(1.2) and the prime-sequence conjecture. \url{https://arxiv.org/abs/2607.08712}.
\end{description}
\label{end:1864}
\end{document}
